% modules/kls/glossary.tex --- back matter, after the synthesis and before the appendices.
%
% Recurring terminology and acronyms only. Every entry is a pointer to where the object is
% actually defined; nothing here is canonical, and no entry is a claim environment, so this
% file carries no ledger node. First use in the text still expands each acronym.
\documentclass[../../main.tex]{subfiles}
\begin{document}

\section{Glossary of recurring terms}
\label{sec:glossary}

Each entry says what the object is and where it is fixed. Where a normalization is at stake
the pointer is to the \texttt{definition} that fixes it, because an exponent or a constant
quoted without its convention is not a statement.

\begin{description}[leftmargin=0pt,style=nextline,itemsep=0.35em]
  \item[Routes E, S, C, F.] The four live routes of this repository, in the order they are
  developed. \textbf{E} follows a fixed candidate bottleneck cut under stochastic
  localization (Section~\ref{sec:introduction}); \textbf{S} follows a fixed first
  eigenfunction instead of a cut (Section~\ref{sec:spectral-route}); \textbf{C} is
  deterministic and works in moment-map coordinates, with no localization at all
  (Section~\ref{sec:moment-map-cmh}); \textbf{F} replaces Euclidean directions by one
  test-independent isotropic frame of conditional line resamplings
  (Section~\ref{sec:conditional-fiber-frame}). Each is an approach family of
  \texttt{research/program/portfolio.yaml}, which is search state and carries no truth
  value.

  \item[Subroutes E--A and E--B.] The two branches of Route~E: E--A is the all-cut Carleson
  estimate (Section~\ref{sec:carleson}), E--B the weighted near-Cheeger package
  (Section~\ref{sec:stein}). See Section~\ref{subsec:two-subroutes}.

  \item[CMH --- the canonical moment-Hessian constant $\CMH$.] The deterministic quantity
  Route~C is built on, fixed with its operator data in Definition~\ref{def:cmh}. It
  dominates the affine Poincar\'e constant $\CPaff$ by
  Theorem~\ref{thm:cmh-implies-affine-poincare}, and it is \emph{not} a reformulation of
  KLS: Proposition~\ref{prop:cmh-hodge} shows it additionally charges a solenoidal excess.

  \item[QCTS --- the quadratic-chaos two-tail statement.] The static input isolated in
  Definition~\ref{def:qcts}, together with the two-tail obstruction
  (Obstruction~\ref{obs:two-tail}) that limits what it can supply.

  \item[Two-colour covariance and two-colour notation.] The bookkeeping that separates the
  contribution of the cut being followed from that of everything else, so that a source term
  can be told from a damping term. Fixed in Section~\ref{subsec:two-color-notation}, used
  throughout Parts~III.

  \item[Stein source, Riccati source and damping.] In the scalar Riccati equation
  $\dd r_t=\dd M_t+(S_t-D_t)\dd t$ of Theorem~\ref{thm:scalar-riccati}, $S_t$ is the only
  positive source and $D_t$ the coercive damping. ``Absorbing the source into the damping''
  is what every fixed-cut argument is trying to do; Appendix~\ref{sec:stein-dictionary}
  gives the Stein representation of the source.

  \item[Interface functional $\Xi_T(\mu)$.] The bridge object of the bootstrap comparison,
  defined at \eqref{eq:interface-def} and evaluated in Section~\ref{sec:bootstrap}. It is
  this repository's version of ``replace $\log$-trace-exp by an effective rank''; the crude
  evaluation is proved insufficient there.

  \item[Gate zero, and sharp gate zero.] The cheapest necessary consequence of the Route~C
  endpoint, Conjecture~\ref{conj:gate-zero}: in isotropic position it asks
  $\lmax(\E H^2)\le4$ where the literature supplies only the trace bound
  $\Tr(\E H^2)\le2n$. That trace bound is attained, so the natural operator statement is the
  \emph{sharp} form $\E H^2\preceq2\Id$ (Conjecture~\ref{conj:gate-zero-sharp}), which implies
  the other. Only the constant~$4$ is what $\mathrm{CMH}(4)$ requires, so the two are not
  interchangeable as falsification targets. Section~\ref{subsec:gate-zero}.

  \item[Exponential cone measure $\bar\mu_{K,\beta}$.] The law with density
  $\propto x_1^{\beta-n}e^{-x_1}$ on the cone over a centered convex body $K$, centered;
  Definition~\ref{def:exponential-cone}. Its moment map is explicit in terms of the base's
  (Proposition~\ref{prop:cone-moment-map}), and at $\beta=n$ it saturates sharp gate zero along
  its axis, which makes the family the first non-product equality set of that statement.
  Section~\ref{subsec:cmh-cones}.

  \item[Covariance spike.] The proved obstruction of
  Proposition~\ref{prop:covariance-spike}: products of centered exponentials are
  dimension-free by tensorization, yet their conditional covariance spikes. It is why no
  uniform operator-norm bound on $\norm{A_t}_\op$ can exist, and it is the single sharpest
  constraint on what a correct argument may look like.

  \item[Applicability-blocked.] A node that is \emph{proved} but whose \texttt{assumes} list
  still holds an open antecedent. It is a theorem, and it is not progress on
  Conjecture~\ref{conj:kls}. The distinction is enforced by the ledger, not by prose; see
  \texttt{CLAUDE.md} constraint~8.

  \item[Candidate (\texttt{cand:}).] A statement proposed in a checkpoint but not yet stable
  enough to be a ledger node. It has no manuscript anchor, no status and no certification,
  and it appears in this manuscript only by name.
\end{description}

\end{document}
